\documentclass[12pt]{article}
\usepackage{mathptmx}
\usepackage[margin=0.75in]{geometry}
\usepackage{natbib}
\usepackage{graphicx} % Required for inserting images

\begin{document}

In analyses of genetic predictors for health outcomes such as disease status, we often employ population-based designs such as the case-control study.  A flaw of this design, however, is its vulnerability to population stratification, which motivates the use of family-based designs, such as case-parents studies where cases and their parents are genotyped (forming a "trio"), allowing one's parents to be considered as controls.  Family-based studies still suffer from unaddressed confounders which may impact potential inferences.  One such confounder is maternal genetic factors, which are relevant when studying diseases known to be dependent on one's gestational environment such as spina bifida \citep{Hobb14}.  Maternal gene-environment interactions, such as the extent smoking impacts a mother's child in her womb being dependent on the mother's genotype, are another potential confounder. We also have mating asymmetry (MA), where mothers and fathers do not have the same distribution of genetic factors among the population, or $\Pr(M=g_1, F=g_2) \not = \Pr(M=g_2, F=g_1)$ where $M$ and $F$ are the genotypes of the mother and father respectively, and $g_1$ and $g_2$ are two distinct genetic factors.  This in turn can cause us to identify false positives such as maternal effects being present should we assume the two distributions are the same when it is really due to population structures.  Although this subject has not been thoroughly studied, \cite{Bour+11} have shown that MA can be observed in human populations, to the extent of creating significant effects in models that do not consider MA.  Lastly, we have parent-of-origin effects (imprinting), where a child's outcome is dependent not only on their genotype but also on the parent who contributed the genetic factor of interest.  In terms of a probability model, we would include $I_m$ and $I_p$ as imprinting parameters, where the probability of disease is multiplied by $I_m$ if they received a high-risk genetic factor maternally, and multiplied by $I_p$ if received paternally \citep{Ains11}.

\cite{Wein98} proposed a log-linear model for case-family data allowing one to draw inferences such as the relative risk of a disease given some genetic factor.  Although no official software for their model exists, one could have used a now deprecated software called LEM, which also had methods for missing data \citep{Verm97}.  Other software for modelling case-family data includes \cite{Howe12}'s EMIM, which employed a multinomial model, and \cite{Gjer19}'s Haplin, which employed a log-linear model.  As models for case-family data, all of them could account for maternal genetic factors and imprinting, however, a limitation of this family-based data was that no inferences about environmental effects could be obtained.  In response to this weakness of the family-based study, \cite{Wein05} introduced a hybrid design, where trio data of cases and their biological parents could be compared to parents of controls to draw inferences using a log-linear model.  With this model, one could include environmental effects such as smoking while still considering all of the previously mentioned cofactors.  Unfortunately, no software methods for this model are currently available, and so users must implement it themselves to perform this analysis.

In our research project, we will create a freely available R package allowing anyone to analyze trio data with the hybrid model.  This includes the basic code to fit the model as a generalized linear model (GLM) object in R, which will allow users to construct confidence intervals and perform significance tests on covariates using functions in base R.  The model will be designed to be robust, allowing users to input data with or without controls and account for missing data using the Expectation-Maximization (EM) algorithm.  We will also conduct several simulation studies to assess the model's accuracy and power in simulated datasets, which could help elucidate the strengths and weaknesses of the model.  As a GLM, the model in our package can have a wide variety of cofactors added to it by advanced users.  By default, our package will consider the common confounders we described before as covariates in the model, with users having the option to disable them.

With this package, our goal is to allow anyone with the appropriate trio data to perform an extensive analysis of the potential contributors to disease.  We expect that the inferences drawn from the model will help inform public health recommendations and for screening high-risk patients.

\newpage
\bibliographystyle{chicago}
\bibliography{Bibliography}

\end{document}
